\chapter{性能评测方法}
\label{ch:benchmarking}

\begin{keybox}
\textbf{本章先回答“怎样比较才算公平”。}并行求解器基准测试的第一步不是计时，而是确认两种实现实际上求解了同一个离散问题、使用了可比的算法参数，并且达到相同的数值精度。
\end{keybox}

\section{先明确 Benchmark 在回答什么问题}

“公平”不是一个单一规则。对求解器性能，至少有两类实验，必须在开始计时以前先声明。

\subsection{受控实现实验：隔离执行后端}

如果问题是“同一个数值算法放到 CPU 与 GPU 后，执行后端本身带来多少差异”，就应该尽量冻结数学路径：
\begin{itemize}
\item PDE、系数、边界条件；
\item 网格与 AMR hierarchy；
\item smoother 类型与 $\nu_1,\nu_2$；
\item cycle 类型；
\item bottom solver；
\item stopping tolerance 与 norm 定义；
\item 初值与浮点精度；
\item hierarchy/setup 是否重用。
\end{itemize}

例如 CUDA 版本使用 weighted Jacobi $2+2$ sweep，而 CPU 版本使用 RBGS $1+1$ sweep，即使两者都得到正确解，也不能把 wall-time 差异解释为“GPU kernel 比 CPU loop 快多少”。此时算法与执行后端已经同时改变。

一张最小受控 benchmark 卡可以固定为：三维常系数 Poisson、同一 $256^3$ 网格、同一 Dirichlet 边界、weighted Jacobi $2+2$ sweep、同一 V-cycle、相同初值和相同相对残差容差，并记录
\begin{equation}
\{N_{cycle},\ T_{setup},\ T_{cycle},\ T_{solve},\ \|r_{final}\|\}.
\end{equation}
如果两端残差历史或 $N_{cycle}$ 已经不同，就应先定位数值路径差异。

\subsection{平台最优实验：比较真实 Time-to-Solution}

另一类问题是“在相同物理问题和精度要求下，CPU/GPU/不同库各自经过合理调优后，谁能更快得到合格答案”。这时不应强迫所有平台使用完全相同的 smoother、tile、coarse strategy 或 backend policy。

这里真正必须固定的是：
\begin{equation}
\boxed{
\text{problem definition}
+
\text{accuracy/robustness target}
+
\text{measurement boundary}.
}
\end{equation}

不同实现可以使用各自合理的算法配置，但必须同时报告：
\begin{equation}
\boxed{
\text{convergence behavior}
\quad+\quad
\text{cost per iteration/cycle}
\quad+\quad
\text{time-to-solution}.
}
\end{equation}

因此本书后面的性能比较应明确标注属于哪一种实验。受控实验用于因果归因，平台最优实验用于工程决策；二者不能用同一套“公平性”规则混在一起。

\section{基准问题设计}

一个最小但有诊断价值的 benchmark suite 应包括：
\begin{enumerate}
\item 2D constant-coefficient Poisson，规则矩形；
\item 3D constant-coefficient Poisson；
\item variable coefficient diffusion；
\item 强各向异性问题，测试 semicoarsening/smoother；
\item 多 patch block-structured geometry；
\item AMR coarse-fine interface；
\item 纯 Neumann/periodic null-space case；
\item 若实际应用有 EB，再加入 cut-cell case。
\end{enumerate}

每个 case 都应有可独立验证的正确性路径，但必须区分两类 manufactured solution。

\paragraph{连续 manufactured solution.}
先选连续函数 $u^\star(x)$，再由连续 PDE 计算 $f$ 和边界数据，最后离散并求解。此时误差
\begin{equation}
E_2=\frac{\|u_h-u^\star_h\|_2}{\|u^\star_h\|_2},
\qquad
E_\infty=\|u_h-u^\star_h\|_\infty
\end{equation}
同时包含离散误差与迭代误差。随 $h$ 加密测实验收敛阶，才能验证离散格式、边界闭合和求解器共同组成的整条链。

\paragraph{离散 manufactured solution.}
直接在网格上指定 $u_h^\star$，再用\textbf{待测离散算子本身}生成
\begin{equation}
b_h=A_hu_h^\star.
\end{equation}
这种测试可以非常敏感地检查 operator apply、边界处理、ghost、transfer 与 solver 是否自洽，但它不能独立证明连续 PDE 的离散阶正确，因为 $b_h$ 已经由同一个离散算子生成。

因此：
\begin{equation}
\boxed{
\text{连续 MMS 用于验证离散收敛，}
\qquad
\text{离散 MMS 用于验证实现一致性。}
}
\end{equation}

\section{强缩放与弱缩放}
\label{sec:scaling}

\textbf{强缩放（strong scaling）}固定全局问题规模，只增加并行资源。实际大问题常常无法在单 rank 或单 GPU 上运行，因此不应把基准强行写死为 $P=1$。选一个能够运行且定义明确的基准资源数 $P_0$，定义
\begin{equation}
S(P;P_0)=\frac{T(P_0)}{T(P)},
\qquad
\eta_{strong}(P;P_0)=\frac{S(P;P_0)}{P/P_0}.
\label{eq:strong-scaling}
\end{equation}
它回答的是：固定同一个全局离散系统以后，增加资源还能把 wall time 降低多少？\chapref{ch:mpi}的式~\eqref{eq:mpi-surface-volume-ratio}已经从 surface/volume 与 latency 解释了为什么这种效率最终会下降。

\textbf{弱缩放（weak scaling）}则保持每个 rank/GPU 的局部工作量近似不变，让全局问题随资源数一起增大。以同样的基准资源数 $P_0$ 定义
\begin{equation}
\eta_{weak}(P;P_0)=\frac{T(P_0)}{T(P)}.
\label{eq:weak-scaling}
\end{equation}
若算法和机器网络理想可扩展，$T(P)$ 应接近常数，因而 $\eta_{weak}$ 接近 1。它更适合观察网络拓扑、global reductions、coarse-grid process reduction 以及 hierarchy setup 是否能跟随系统规模增长。

\begin{keybox}
强缩放和弱缩放不是两个不同的“画图方式”，而是在回答两个不同问题：前者问“同一个问题能否更快”，后者问“机器变大时能否保持单位资源效率”。
\end{keybox}

\section{重复测量与统计}

并行性能具有运行间波动，因此单次 wall time 不能代表稳定性能。一个基本流程是：先做若干次 warm-up（正式计时前的预运行），使 JIT（just-in-time compilation，运行时首次生成/优化代码）、缓存、内存池（预先保留并复用的分配空间）和 GPU context（设备运行时上下文）等一次性成本进入稳定状态；随后做多次独立计时，并报告样本分布而不是只给一个平均值。

若测得时间样本 $t_1,\ldots,t_m$，至少建议报告中位数（median）以及一个离散程度指标，例如四分位距（IQR）或最小/最大值。对存在明显系统噪声的 HPC 环境，中位数通常比算术平均更不容易被少量异常慢样本拖动。

还必须明确计时边界。Setup、第一次数据迁移、一次 V-cycle、完整 solve、输出与 checkpoint 不能混在同一个数字里。后文\chapref{ch:diagnosis}的 profiler 诊断只有在计时边界稳定时才有意义。

\paragraph{Wall time、互斥阶段时间与可重叠事件时间。}
异步 CPU/GPU/MPI 程序至少要区分三种时间语义。第一，\textbf{wall time} 是从阶段开始到结束的实际经过时间，本质上由执行 DAG 的关键路径决定。第二，\textbf{exclusive phase time} 把每一段 wall-clock 时间唯一归入某个阶段，因此可以满足
\begin{equation}
T_{wall}=\sum_k T_k^{\mathrm{exclusive}}.
\label{eq:benchmark-exclusive-time}
\end{equation}
第三，profiler 常直接报告 kernel、memory copy、MPI request 等\textbf{event duration}；这些事件可以彼此重叠，因此一般只有
\begin{equation}
T_{wall}\neq\sum_j T_j^{\mathrm{event}}.
\label{eq:benchmark-overlap-event-time}
\end{equation}
例如 halo 通信与 interior residual 并发时，两者 event duration 都是真实成本，但其重叠部分只在 wall time 中出现一次。报告逐层 breakdown 时必须注明采用的是 exclusive accounting 还是 raw event duration；只有前者可以无条件求和。

\section{逐层性能计数}

建议对每层记录：
\begin{equation}
\{N_\ell,\#boxes_\ell,\#ranks_\ell,T_{smooth},T_{res},T_R,T_P,T_{halo},T_{bottom}\}.
\end{equation}
若这些 $T$ 用于可加和的 phase breakdown，默认应采用式~\eqref{eq:benchmark-exclusive-time}的 exclusive accounting；若来自 profiler 的 raw event duration，则应显式标注为 event time，并允许总和大于 wall time。
对显式稀疏/AMG hierarchy，还应记录 $\mathrm{nnz}(A_\ell)$，并用\chapref{ch:software}的式~\eqref{eq:grid-complexity}与式~\eqref{eq:operator-complexity}报告 grid/operator complexity。否则两个 solve 时间相近的配置，可能具有完全不同的 setup 内存和粗层稠密化成本。
GPU 再增加：
\begin{equation}
\{\text{kernel duration},\text{achieved bandwidth},\text{occupancy},\text{L2 hit},\text{launch count}\}.
\end{equation}
CPU/OpenMP 增加：
\begin{equation}
\{\text{memory bandwidth},\text{LLC misses},\text{NUMA remote traffic},\text{barrier time}\}.
\end{equation}
MPI 增加：
\begin{equation}
\{\text{message count},\text{bytes},\text{wait time},\text{Allreduce time}\}.
\end{equation}

如果只看总时间，你只能知道“慢”；逐层数据才能回答“从第几层开始变慢，以及为什么”。

\section{端到端求解成本}

\subsection{收敛因子与端到端求解时间}

测每 V-cycle residual reduction：
\begin{equation}
\rho_k=\frac{\|r^{(k+1)}\|}{\|r^{(k)}\|}.
\end{equation}
稳定区间可用几何平均估计 asymptotic convergence factor：
\begin{equation}
\bar\rho=\biggl(\frac{\|r^{(m)}\|}{\|r^{(0)}\|}\biggr)^{1/m}.
\end{equation}
然后同时报告：
\begin{equation}
T_{cycle},\qquad \bar\rho,\qquad T_{solve}.
\end{equation}

例如 smoother A 使 $\rho=0.08$ 但 cycle 时间 \SI{10}{ms}；smoother B 使 $\rho=0.2$ 但 cycle 只要 \SI{3}{ms}。把几何收敛模型连续化估计：
\begin{equation}
N_A^\star=\frac{\ln10^{-8}}{\ln0.08}\approx7.3,
\qquad
N_B^\star=\frac{\ln10^{-8}}{\ln0.2}\approx11.4.
\end{equation}
实际 cycle 数必须取满足容差的最小整数，因此
\begin{equation}
N_A=8,
\qquad
N_B=12.
\end{equation}
对应总时间约为 \SI{80}{ms} 与 \SI{36}{ms}。A 的单 cycle 收敛明显更强，但 B 的 time-to-solution 仍然更短。

\subsection{构建与求解成本}

GMG hierarchy 通常便宜，但 AMR regrid 后也可能需要重建；AMG setup 则可能很昂贵。总成本应写成
\begin{equation}
T_{total}=T_{setup}+N_{reuse}T_{solve}.
\end{equation}
当同一矩阵重复求解很多次，昂贵 setup 可以摊薄；若每个时间步矩阵结构都改变，setup 性能可能比 solve 更重要。

\section{把抽象阶段映射到真实源码}

\chapref{ch:software}已经说明不同软件把相同数学阶段放在不同对象中。Benchmark 若只记录一个笼统的 \texttt{solve()} 时间，就无法解释库间差异。下面把本章的 phase accounting 映射到第 14 章固定源码快照中的真实入口。

\begin{center}
\small
\begin{tabularx}{0.98\textwidth}{>{\bfseries}p{2.0cm}p{3.7cm}X}
\toprule
实现 & 源码入口/事件 & Benchmark 中应归属的阶段\\
\midrule
AMReX MLMG &
\texttt{MLMGT::solve}、\texttt{oneIter}、\texttt{mgVcycle}、\texttt{actualBottomSolve} &
外层 solve、单次 multigrid iteration、V-cycle、bottom solve；\texttt{BL\_PROFILE} 已在这些控制流中提供可追踪区域\\
hypre BoomerAMG &
\texttt{hypre\_BoomerAMGSetup}、\texttt{hypre\_BoomerAMGCycle}、ParCSR Matvec/RAP &
setup 与 cycle 必须分离；cycle 内继续拆 relaxation、residual SpMV、restriction/prolongation 与 coarse solve\\
PETSc PCMG/GAMG &
\texttt{PCSetUp\_GAMG}、\texttt{PCMGMCycle\_Private}、每层 \texttt{KSP}、\texttt{MatRestrict}/\texttt{MatInterpolateAdd} &
GAMG hierarchy setup、MG apply、level smoother、transfer；源码中已有 \texttt{PetscLogEventBegin} 包围 smoother、restriction 等阶段\\
\bottomrule
\end{tabularx}
\end{center}

这张映射表的作用不是要求三个软件输出完全相同的 profiler 标签，而是建立一个共同的\textbf{语义计时层}。原始 profiler event 应先映射到
\begin{equation}
\boxed{
T_{setup},
T_{smooth},
T_{residual},
T_R,
T_P,
T_{halo},
T_{redistribute},
T_{bottom},
T_{reduce}
}
\label{eq:benchmark-semantic-bins}
\end{equation}
再进行跨实现比较。

这样可以避免一个常见错误：AMReX 的 \texttt{mgVcycle}、hypre 的 \texttt{BoomerAMGCycle}、PETSc 的 \texttt{PCApply} 虽然都是“一个大函数”，但它们内部包含的 setup、level object、smoother 与通信责任并不相同，直接比较顶层函数耗时不能说明具体实现机制。

\section{性能剖析层级}
\label{sec:setup-solve}

\begin{verbatim}
solve
  setup / hierarchy rebuild
  vcycle
    level 0
      fill boundary / halo
      pre-smooth
      residual
      restrict
    level 1
      ...
    bottom solve
    level 1
      prolong + correction
      post-smooth
    level 0
      prolong + correction
      post-smooth
  residual norm / convergence test
\end{verbatim}

这棵树应该在 CPU、MPI、GPU 版本中保持相同语义。只有这样，profile 才能跨实现比较。

\section{本章小结}
本章首先区分了两类不同实验：受控实现实验要求冻结数值路径，用于隔离执行后端；平台最优实验允许各实现采用合理配置，用于比较真实 time-to-solution。正确性验证又进一步区分连续 MMS 与离散 MMS：前者验证离散收敛，后者验证离散实现自洽。

在此基础上，性能评测才进入 warm-up、重复采样、strong/weak scaling、逐层计数、收敛因子与 setup amortization。最终所有 profiler event 都应映射回稳定的语义阶段，而不是直接比较库内部的大函数名称。所有性能结论必须绑定到明确的问题定义、精度目标、计时边界、资源规模和样本统计。

\section*{习题}

\subsection*{基准量与统计推导}
\begin{enumerate}[label=\arabic*.]
\item 某 benchmark 运行时间为 $[10.2,10.0,9.9,10.1,35.0]$ ms。先排序；采用 Tukey 约定（奇数样本的总体中位数不再计入上下半组）计算均值、中位数、$Q_1,Q_3$ 与 IQR，并判断哪一个统计量更适合报告典型性能。
\item 强缩放基准在 $P=4,8,16,32$ 时耗时为 $80,43,24,16$ s。以 $P_0=4$ 为基准计算 speedup 与 parallel efficiency。
\item 弱缩放中每 rank 固定 $10^6$ unknowns，$P=1,8,64$ 时耗时为 $1.0,1.08,1.35$ s。计算相对弱缩放效率。
\item 一个 solver setup=2 s、每次 solve=0.1 s；另一个 setup=0.2 s、每次 solve=0.14 s。求两者总时间相等时的 reuse 次数，并判断重复 1、20、100 次时谁更快。
\item smoother A 每 cycle 8 ms、convergence factor 0.15；B 每 cycle 5 ms、factor 0.30。假设初始 residual 为 1，计算达到 $10^{-8}$ 所需 cycles 与总时间。

\item 从 speedup 定义推导以任意基准资源数 $P_0$ 定义的 strong-scaling efficiency，说明为什么不必假设单 rank 能运行。
\item 推导 setup/solve 总成本 $T=T_{setup}+N_{reuse}T_{solve}$ 的 break-even reuse 次数公式。
\item 对几何收敛 $\|r_k\|\approx q^k\|r_0\|$ 推导达到 tolerance $\varepsilon$ 所需 cycle 数。
\item 说明为什么只有在 PDE、离散、停止准则和初值等数值路径等价时，性能比较才有可解释性；构造一个“更快只是因为少做了迭代”的反例。
\end{enumerate}

\subsection*{基准实验设计}
\begin{enumerate}[label=\arabic*.]
\item 写一个 benchmark harness：warm-up 后重复至少 15 次，输出 median、IQR、min/max，并保存原始样本。
\item 分别实现连续 MMS 与离散 MMS：前者自动对多组 $h$ 计算误差和实验阶；后者用 $b_h=A_hu_h^\star$ 检查 operator/boundary/solver 一致性。比较两种测试各自能发现和不能发现的问题。
\item 对同一 V-cycle 自动生成 strong/weak scaling 表和图，同时记录每 level local problem size。
\item 实现 setup/solve 分离计时，并模拟不同 $N_{reuse}$ 计算摊销后的总成本。
\end{enumerate}

\subsection*{证据解释与公平性诊断}
\begin{enumerate}[label=\arabic*.]
\item 分别设计“受控 CPU vs GPU 后端实验”和“平台最优 time-to-solution 实验”两张 benchmark card，指出哪些算法参数前者必须固定、后者可以平台特化。
\item 为什么只报告“最快一次”容易误导？结合 warm-up、JIT、cache、系统噪声说明。
\item 设计一个 release-blocking benchmark 集：至少包含常系数、变系数、AMR、null-space 与并行缩放问题，并说明每类在验证什么。
\end{enumerate}
